% Service statement -- Benjamin Peirce Fellowship, Harvard University

During my time as a PhD student, I have learned much
in seminars, in reading groups, and in conversations with other early-career mathematicians.
Similarly, as an undergraduate and Master's student, I profited enormously from the mentorship of PhD students.
I see contributing to these communities as an important part of being a mathematician,
and wish to give back in a similar way.
During my PhD at Cambridge, I have done this mainly by organizing seminars and reading groups,
and by helping incoming students make the transition to graduate-level mathematics.

From October 2024 to June 2026, I co-organized the Junior Geometry Seminar at the University of Cambridge.
The seminar gives PhD students and postdocs in geometry and topology a place to present their work,
often for the first time outside their own research group.
The speakers are a mix of internal speakers and early-career researchers from other institutions, mainly from those in the UK.
This allows networking and has sparked more than one collaboration.
Since the audience comes from all areas of geometry, we encouraged speakers to make their talks accessible to non-experts,
which also gave the speakers valuable practice in communicating their research to a broad audience.
Giving a talk there early-in my PhD was very encouraging, and this experience motivated me to join the organizers.

From December 2023 to June 2025, I also co-organized the A-side Seminar,
a learning seminar at Cambridge on symplectic topology (named for the A-side of homological mirror symmetry), with Alexia Corradini.
Each week one of the PhD students or postdocs in symplectic geometry gave a talk
on a topic of general interest. However, it was never a closed internal seminar.
From the start, we encouraged interested students, especially from Part III (Cambridge's one-year Master's course),
to join the audience.
This came to fruition, as one of the Part III students that joined our seminar back then has since joined our
group as a PhD student and now shares an office with me.

Students arrive in Part III from many different universities and countries,
and with very different preparation.
In the summer of 2024, I designed and taught the algebraic topology course of the Part III Preparatory Workshop,
consisting of two sessions and a problem sheet.
My aim was to help students whose previous studies had not covered algebraic topology in depth start the year on an equal footing with their peers,
and to give everyone a chance to ask questions in a low-stakes setting before the term began.
Having moved between universities from Berlin to Zurich, and from Zurich to Cambridge myself,
I remember well how much a good start in a new program matters.

I am also interested in building connections between different fields.
With Anand Srinivasan, a PhD student in the fluid dynamics group, I organized an informal reading group in 2025. Its topic was
the work of Cardona, Miranda, Peralta-Salas, and Presas showing that the
Euler equations of an ideal fluid admit solutions that simulate an
arbitrary Turing machine. The question spans fluid mechanics,
partial differential equations, contact geometry (the odd-dimensional
analogue of symplectic geometry), and computer science, and our members
reflected this, ranging from pure to highly applied mathematicians.

At Harvard, I would be glad to contribute to the department in similar ways.
For graduate students and postdocs, I would like to help with the organization of seminars in geometry and topology,
and also a reading group bridging fields, similar to the one I initiated at Cambridge.
For undergraduates, I would be happy to give a talk at the Open Neighborhood Seminar,
to encourage students working on projects with me to present their work at Math Table,
and to supervise reading and research projects.
Parts of my research have accessible computational aspects, which are well suited for such projects.

I have benefited greatly from the mathematical communities I have been part of in Berlin, Zurich, and Cambridge,
and I look forward to contributing to the community at Harvard.
